\chapter{Introducción}

% --------------------------------------------------------
\section{Contexto y motivación}
% --------------------------------------------------------
% TODO: ~500 palabras
% Puntos:
%   - El EMBI de Ecuador como indicador de riesgo soberano: qué mide, por qué importa
%   - Ecuador: economía dolarizada, historial de defaults (1999, 2008 Correa, 2020 pandemia)
%   - Costo económico: 1 pb en el EMBI representa decenas de millones USD en servicio de deuda
%   - Brecha: los modelos econométricos tradicionales (ARIMA) no capturan shocks exógenos ni información mediática
%   - Oportunidad: disponibilidad de datos GDELT v2 (noticias globales) + ML moderno

% --------------------------------------------------------
\section{Pregunta de investigación}
% --------------------------------------------------------
% TODO: ~200 palabras
% Pregunta principal: ¿Puede un modelo de aprendizaje automático entrenado con
%   variables macroeconómicas, financieras externas e información mediática (NLP)
%   predecir el EMBI de Ecuador con mayor precisión que el baseline ARIMA?
%
% Preguntas secundarias:
%   1. ¿Qué grupo de variables (AR, macro, NLP) contribuye más a la predicción?
%   2. ¿Anticipa el volumen de cobertura mediática los movimientos del EMBI (causalidad de Granger)?
%   3. ¿Existen quiebres estructurales que cambien qué variables determinan el EMBI?
%   4. ¿Mejora la compresión no lineal (autoencoder) la representación NLP?

% --------------------------------------------------------
\section{Objetivos}
% --------------------------------------------------------

\subsection{Objetivo general}
% TODO

\subsection{Objetivos específicos}
% TODO (4-5 objetivos alineados con las preguntas secundarias)

% --------------------------------------------------------
\section{Hipótesis}
% --------------------------------------------------------
% TODO:
%   H1: El XGBoost con variables AR + NLP supera al ARIMA en RMSE fuera de muestra
%   H2: El NLP aporta mejora incremental significativa sobre el modelo solo-AR
%   H3: No existe causalidad de Granger NLP → EMBI con rezago diario
%   H4: Existen quiebres estructurales estadísticamente significativos en 2020 y 2022

% --------------------------------------------------------
\section{Contribuciones}
% --------------------------------------------------------
% TODO: ~300 palabras
% 1. Primera aplicación sistemática de GDELT v2 para predicción del EMBI ecuatoriano
% 2. Experimento de ablación que aísla la contribución del NLP vs macro vs AR
% 3. Análisis de regímenes (Chow + CUSUM + SHAP) que muestra heterogeneidad estructural
% 4. Evidencia empírica sobre exogeneidad del EMBI ecuatoriano en períodos normales vs crisis

% --------------------------------------------------------
\section{Estructura del documento}
% --------------------------------------------------------
% TODO: párrafo breve describiendo cada capítulo
